\documentclass[10pt,a4paper]{amsart}
\usepackage[margin=19mm]{geometry}
\usepackage{amsmath,amssymb,mathtools}
\usepackage[scr=rsfs,cal=euler]{mathalfa}
\usepackage{xfrac}
\usepackage[unicode,hidelinks,bookmarks=false]{hyperref}
\newcommand{\ie}{i.e.\ }
\newcommand{\cf}{cf.\ }
\newcommand{\resp}{respectively\ }
\newcommand{\Subsection}{\subsection}
\providecommand{\nobreakdash}{\nobreak\hbox{-}}
\setlength{\emergencystretch}{2em}
\multlinegap0pt
\usepackage{amssymb}
\usepackage{mathtools}
\usepackage{tikz}
\usepackage[shortlabels]{enumitem}
\usepackage{xcolor}
\newtheorem{lem}{Lemma}
\newtheorem{thm}{Theorem}
\usepackage{cleveref}
\crefname{lem}{Lemma}{Lemmas}
\crefname{thm}{Theorem}{Theorems}
\title{Independence Polynomials of Graphs}
\author{Takayuki Hibi, Selvi Kara, Dalena Vien}
\date{Source: arXiv:2603.16695v1 (CC BY 4.0)}
\begin{document}
\maketitle

\noindent\emph{Source and licence.} Independence Polynomials of Graphs. Version arXiv:2603.16695v1 (CC BY 4.0), distributed by arXiv under the Creative Commons Attribution 4.0 International licence, http://creativecommons.org/licenses/by/4.0/, as recorded in the arXiv licence metadata for this exact version and shown on the abstract page https://arxiv.org/abs/2603.16695v1. Full licence notice: \texttt{licenses/arxiv-2603.16695-CC-BY-4.0.txt}.\par\smallskip

\noindent\textit{Editorial scope.} Counted unit: the article's thm prop:big\_star\_symmetry (source0.tex lines 637--639). The article's complete original proof of it is delivered with it (source0.tex lines 641--794, one \texttt{proof} environment of 154 lines). No mathematics is written, repaired or re-derived: the statement and the proof are the article's own lines, in the article's own notation, and no retained line is substituted or re-flowed.  Retained uncounted as the source-local context the proof needs: lines 578--586.  Nothing was omitted from any retained span, so the package carries no omission record.  No retained line contains a citation, so the wrapper carries no bibliography and no bibliography entry.  No retained line contains a figure environment, an image-inclusion command or drawing code, so no image is delivered and the package declares no asset dependency.  Editorial formatting only, in addition: an AMS \texttt{amsart} preamble that restates the article's own declarations and macros used by the retained lines, namely the environments \texttt{lem}, \texttt{thm} on separate counters, together with the standard packages those lines need.  The retained environments print in this document as Lemma 1, Theorem 1 in order of appearance, whereas the article prints the same items with the numbering of its own full document.  The complete native source of the article as distributed in the arXiv e-print for this exact version is bundled unchanged as \texttt{sources/196569/source0.tex}.
\begin{lem}\label{lem:alpha_big_star}
Let $G(n_1,\ldots,n_q)$ be a big star graph. Then
\[
\alpha\bigl(G(n_1,\ldots,n_q)\bigr)
=
\sum_{i=1}^q \left\lfloor n_i/2\right\rfloor+\max\{1,r\}
\]
where  \(r=|\{ i : n_i \text{ is odd} \}|\).
\end{lem}

\begin{thm}\label{prop:big_star_symmetry}
Let \(G=G(n_1,\ldots,n_q)\) be a big star. Then \(P_G(x)\) is symmetric if and only if $G\cong G(1,1,5)$.
\end{thm}

\begin{proof}
Let $\alpha:=\alpha(G)$. If \(P_G(x)\) is symmetric, then \(g_{\alpha}=g_0=1\). This means \(G\) has a unique
maximum independent set.  


Recall that  \(r=|\{ i : n_i \text{ is odd} \}|\).  If \(r=1\), then there is a maximum independent set containing the center and also one
avoiding the center, contradiction. If \(2\le r\le q-1\), then every maximum independent
set avoids the center, and each even arm \(P_{2a}\) contributes at least two maximum
choices, again a contradiction. Hence either all arms are even or all arms are odd.

Suppose first that \(n_i=2a_i\) for all \(i\). In this case, $g_1=1+2\sum_{i=1}^q a_i$ and    $\alpha=1+\sum_{i=1}^q a_i$ by \Cref{lem:alpha_big_star}. An independent set of size \(\alpha-1\) either avoids the center or contains it. If it avoids the center, then on the \(i\)th arm $x^{(i)}_1-x^{(i)}_2-\cdots-x^{(i)}_{2a_i}$
we must choose a maximum independent set of \(P_{2a_i}\). Thus each arm $P_{2a_i}$ contribute exactly $a_i$ vertices and the number of independent sets of size $a_i$ in $P_{2a_i}$ is $\binom{2a_i+1-a_i}{a_i}=a_i+1$. Thus these contribute $\prod_{i=1}^q (a_i+1)$.

If it contains the center, then \(x^{(i)}_1\) is forbidden on every arm. So on the \(i\)th arm we may only choose vertices from $x^{(i)}_2-x^{(i)}_3-\cdots-x^{(i)}_{2a_i}\cong P_{2a_i-1}$.
Since the center already contributes one vertex, the arms together must contribute \(\sum_{i=1}^q a_i-1\). Hence exactly one arm, say the \(j\)th, contributes \(a_j-1\) vertices, while every other arm contributes \(a_i\), necessarily the unique maximum independent set $\{x^{(i)}_2,x^{(i)}_4,\ldots,x^{(i)}_{2a_i}\}$. 
On the exceptional \(j\)th arm, the number of independent sets of size \(a_j-1\) in \(P_{2a_j-1}\) is $\binom{(2a_j-1) +1- (a_j-1)}{a_j-1}= \binom{a_j+1}{2}$. 
Thus these contribute $\sum_{j=1}^q \binom{a_j+1}{2}$. Therefore
\[
g_{\alpha-1}
=
\prod_{i=1}^q (a_i+1)+\sum_{i=1}^q \binom{a_i+1}{2}.
\]
Since \(\binom{a_i+1}{2}\ge a_i\) and, because \(q\ge 3\) and each \(a_i\ge 1\), we have $\prod_{i=1}^q (a_i+1)>1+\sum_{i=1}^q a_i$.  Thus, $g_{\alpha-1}>1+2\sum_{i=1}^q a_i=g_1$,
contradicting symmetry. This means all arms are odd. 

Let \(n_i=2a_i+1\) for all \(i\). By \Cref{lem:alpha_big_star}, we have $\alpha
=q+\sum_{i=1}^q a_i$. Now let \(I\) be an independent set containing the center \(x\). Then \(x^{(i)}_1\notin I\) for every \(i\), so on the \(i\)th arm we may choose vertices only from $x^{(i)}_2-x^{(i)}_3-\cdots-x^{(i)}_{2a_i+1}\cong P_{2a_i}$. Hence at most \(a_i\) vertices from that arm. Therefore, since \(q\ge 3\), we get $|I|\le   1+\sum_{i=1}^q a_i
\le \alpha-2$. Thus no independent set of size \(\alpha-1\) contains the center. It follows that every independent set counted by \(g_{\alpha-1}\) avoids the center. On the \(i\)th arm, the maximum possible contribution is \(a_i+1\).  Hence such a set must be maximal on all but exactly one arm, and on the remaining arm, say the \(i\)th, it must have size \(a_i\). For \(j\neq i\), the unique maximum independent set of \(P_{2a_j+1}\) is
\( \{x^{(j)}_1,x^{(j)}_3,\ldots,x^{(j)}_{2a_j+1}\}
\),
while the number of independent sets of size \(a_i\) in \(P_{2a_i+1}\) is
\(\binom{(2a_i+1)+1-a_i}{a_i}=\binom{a_i+2}{2}\).
Summing over the choice of the exceptional arm gives
\[
g_{\alpha-1}=\sum_{i=1}^q \binom{a_i+2}{2}.
\]
Since \(P_G(x)\) is symmetric of degree \(\alpha\), we have  
\[
g_{\alpha-1}= \sum_{i=1}^q \binom{a_i+2}{2}
=
q+1+2\sum_{i=1}^q a_i=g_1.
\]
Since $\binom{a_i+2}{2}=\binom{a_i}{2}+2a_i+1$,
the equality becomes $\sum_{i=1}^q \binom{a_i}{2}=1$. This equality forces exactly one term to be \(1\) and all the others to be \(0\). Hence, after reordering, exactly one \(a_i\) is \(2\), and every other \(a_j\) is \(0\) or \(1\). Since \(n_i=2a_i+1\), this means that exactly one arm has length \(5\), and every other arm has length \(1\) or \(3\). 


If \(q=3\), then the only possibilities are \(G(1,1,5)\), \(G(1,3,5)\), and \(G(3,3,5)\),
and a direct computation shows that only \(G(1,1,5)\) is symmetric.


Now assume \(q\ge 4\). From the previous paragraph, after reordering we may write
\[
G\cong G(\underbrace{1,\ldots,1}_{b},\underbrace{3,\ldots,3}_{t},5)
\]
where $q=b+t+1$ and $\alpha=b+2t+3$. 


Set $S(x):=(1+x)^b(1+3x+x^2)^t$. 
Then, we have 
\[
P_G(x)=S(x)\underbrace{(1+5x+6x^2+x^3)}_{A(x)}+\underbrace{x(1+4x+3x^2)(1+2x)^t}_{B(x)}.
\]
Since \(S(x)\) is symmetric of degree \(b+2t=\alpha-3\), we have $x^{\alpha-3}S(x^{-1})=S(x)$. Then
\[
[x^{\alpha-k}]S(x)A(x)=[x^k]S(x)x^3A(x^{-1}).
\]
Therefore
\[
[x^k]S(x)A(x)-[x^{\alpha-k}]S(x)A(x)
=
[x^k]S(x)\bigl(A(x)-x^3A(x^{-1})\bigr)
\]
where  $A(x)-x^3A(x^{-1})= x(x-1)$.
Hence
\[
[x^k]S(x)A(x)-[x^{\alpha-k}]S(x)A(x)
=
[x^k]\bigl(x(x-1)S(x)\bigr).
\]

Next, $\deg B=t+3=\alpha-(q-1)$.
So if \(k\le q-2\), then \(\alpha-k>\deg B\). Thus $[x^{\alpha-k}]B(x)=0$.
It follows that for every \(k\le q-2\), we have
\[
g_k-g_{\alpha-k}
=
[x^k]\Bigl(B(x)+x(x-1)S(x)\Bigr).
\]
Consider \(k=2\). One can deduce the coefficients of $x^2$ in $B(x)$ and $x(x-1)S(x)$ to conclude that $g_2-g_{\alpha-2}=6-q$.
% Since
% \[
% S(x)=1+(b+3t)x+\cdots
% \]
% and
% \[
% B(x)=x+(4+2t)x^2+\cdots,
% \]
% we get
% \[
% x(x-1)S(x)=-x+(1-b-3t)x^2+\cdots.
% \]
% Therefore
% \[
% g_2-g_{\alpha-2}
% =
% (4+2t)+(1-b-3t)
% =
% 5-b-t
% =
% 6-q.
% \]
If \(P_G(x)\) is symmetric, then $q=6$ and $b+t=5$.

Now consider \(k=3\) and make use of \(b=5-t\) in deducing the coefficient of $x^3$ in  $B(x)$ and $x(x-1)S(x)$. So, we have $g_3-g_{\alpha-3}=t-2$.
% \[
% S(x)=(1+x)^{5-t}(1+3x+x^2)^t
% =
% 1+(5+2t)x+(10+7t+2t^2)x^2+\cdots,
% \]
% while
% \[
% B(x)=x(1+4x+3x^2)(1+2x)^t
% =
% x+(4+2t)x^2+(2t^2+6t+3)x^3+\cdots.
% \]
% Thus
% \[
% x(x-1)S(x)
% =
% -x-(4+2t)x^2+\bigl((5+2t)-(10+7t+2t^2)\bigr)x^3+\cdots,
% \]
% and so
% \[
% g_3-g_{\alpha-3}
% =
% (2t^2+6t+3)+(5+2t)-(10+7t+2t^2)
% =
% t-2.
% \]
Symmetry forces \(t=2\). Hence \(b=3\). Thus the only  possibility with \(q\ge 4\) is $G\cong G(1,1,1,3,3,5)$.
A direct computation gives the following polynomial, which is not symmetric:
\[
P_G(x)
% =
% (1+x)^3(1+3x+x^2)^2(1+5x+6x^2+x^3)
% +
% x(1+4x+3x^2)(1+2x)^2
% \]
% \[
=
1+15x+91x^2+296x^3+577x^4+714x^5+575x^6+296x^7+91x^8+15x^9+x^{10}.
\]
Therefore \(P_G(x)\) is symmetric if and only if \(G\cong G(1,1,5)\).
\end{proof}

\end{document}
